% Guía de mediciones para llevar al laboratorio (Parte 1 del TP6).
\documentclass[10pt]{article}
\usepackage[utf8]{inputenc}
\usepackage[T1]{fontenc}
\usepackage[spanish,es-tabla,es-nodecimaldot]{babel}
\usepackage[a4paper,margin=1.5cm]{geometry}
\usepackage{array}
\usepackage{booktabs}
\usepackage{enumitem}
\usepackage[table]{xcolor}
\usepackage{amsmath}
\usepackage{amssymb}
\usepackage{pdflscape}
\setlist{itemsep=2pt, topsep=2pt, leftmargin=*}
\setlength{\parindent}{0pt}
\pagestyle{empty}
\newcolumntype{L}[1]{>{\raggedright\arraybackslash}p{#1}}
\newcommand{\archivo}[1]{\texttt{\small #1}}

\begin{document}

{\Large\textbf{TP6 Medios -- Guía de mediciones de la Parte 1}}\\[2pt]
Sistemas de Radiocomunicación -- FCEFyN, UNC \hfill Fecha: \underline{\hspace{3cm}} \quad Comisión: \underline{\hspace{2cm}}

\vspace{6pt}
\textbf{Antes de empezar}
\begin{enumerate}
    \item Fotografiar la antena, el analizador, el bloqueador de continua, los cables y el adaptador (\archivo{montaje.jpg}). Anotar marca y modelo de cada equipo.
    \item Conexión: \textbf{antena $\rightarrow$ (adaptador 75/50 $\Omega$) $\rightarrow$ bloqueador de continua $\rightarrow$ analizador}. Solo señales de recepción: \textbf{nunca} conectar un transmisor.
    \item ¿Se usó adaptador 75/50 $\Omega$? \quad $\square$ Sí \quad $\square$ No
    \item En cada captura deben leerse: inicio/fin o centro, \textit{span}, nivel de referencia, RBW, VBW, atenuación y unidad. Si no aparecen en pantalla, anotarlos en la planilla.
    \item Punto de partida: nivel de referencia $-30$ dBm (ajustar para que la señal no salga de pantalla), ganancia 5--10 dB.
\end{enumerate}

\vspace{6pt}
\textbf{Cuadro de mediciones}\\[3pt]
{\small
\renewcommand{\arraystretch}{1.5}
\begin{tabular}{@{}c L{2.2cm} L{2.4cm} L{1.3cm} L{1.4cm} L{1.2cm} L{4.2cm} L{2.0cm}@{}}
    \toprule
    \textbf{\#} & \textbf{Medición} & \textbf{Inicio--fin / centro} & \textbf{Span} & \textbf{RBW} & \textbf{VBW} & \textbf{Qué anotar} & \textbf{Captura} \\
    \midrule
    1 & Banda AM completa & 535--1705 kHz & 1,17 MHz & 9 kHz & 100 kHz & f y nivel de la emisora más fuerte & \archivo{am\_banda} \\
    \rowcolor{gray!10} 2 & Emisora AM en detalle & Centro = f emisora & 20--30 kHz & \textbf{$\leq$ 1 kHz} & 3--10 kHz & f y nivel de portadora, ancho a $-26$ dB, emisoras a $\pm$10 y $\pm$20 kHz & \archivo{am\_emisora} \\
    3 & Banda FM completa & 88--108 MHz & 20 MHz & 100 kHz & 300 kHz & f y nivel de la emisora más fuerte & \archivo{fm\_banda} \\
    \rowcolor{gray!10} 4 & Emisora FM en detalle & Centro = f emisora & 1--2 MHz & 10 kHz & 30--100 kHz & f, nivel, ancho a $-26$ dB (max hold), emisoras vecinas y separación & \archivo{fm\_emisora} \\
    5 & TV analógica (si hay) & Centro del canal & 8 MHz & 30--100 kHz & 300 kHz & Canal; f y nivel de video, color y audio; ancho; canales vecinos & \archivo{tv\_analogica} \\
    \rowcolor{gray!10} 6 & Banda UHF (TDA) & 470--698 MHz & 228 MHz & 100--300 kHz & 1 MHz & Canales visibles: f y nivel aproximado & \archivo{tda\_banda} \\
    7 & Canal ISDB-Tb en detalle & $473 + 6(N-14)$ MHz & 10 MHz & 30--100 kHz & 300 kHz & Nivel de meseta, ancho a $-3$ dB, potencia de canal, forma de la meseta & \archivo{tda\_canal} \\
    \rowcolor{gray!10} 8 & Búsqueda en el TV & -- & -- & -- & -- & Canales y servicios encontrados & foto de la lista \\
    9 & AM con generador (opcional) & Centro = f generador & 5--10 kHz & 100--300 Hz & 1 kHz & Nivel de portadora y de una banda lateral, f del tono & \archivo{am\_generador} \\
    \bottomrule
\end{tabular}}

\newpage
\textbf{Paso a paso}

\vspace{3pt}
\begin{minipage}[t]{0.49\textwidth}
\textbf{1. AM}
\begin{enumerate}
    \item Banda completa (fila 1) y captura.
    \item Marcador en el pico más alto: anotar f y nivel.
    \item Centrar, \textit{span} 20--30 kHz y \textbf{RBW $\leq$ 1 kHz} (con 9 kHz no se ven las bandas laterales).
    \item \textit{Max hold} 10--20 s.
    \item Con marcadores, medir el ancho donde cae 26 dB bajo la portadora.
    \item Ver si hay emisoras a $\pm$10 y $\pm$20 kHz. Captura.
    \item Identificar la emisora (CABFRA / listado ENACOM).
\end{enumerate}

\vspace{4pt}
\textbf{2. FM}
\begin{enumerate}
    \item Banda completa (fila 3) y captura.
    \item Marcador en la más fuerte: f y nivel. Canal: $N = 201 + (f - 88{,}1)/0{,}2$.
    \item Centrar, \textit{span} 1--2 MHz, RBW 10 kHz.
    \item \textit{Max hold} 20--30 s; ancho a $-26$ dB del máximo.
    \item Anotar emisoras vecinas y su separación en kHz. Captura.
\end{enumerate}

\vspace{4pt}
\textbf{3. TV analógica (solo si hay señal)}
\begin{enumerate}
    \item Buscar un canal con portadoras en línea (no meseta plana).
    \item Centrar, \textit{span} 8 MHz.
    \item Marcadores: video ($\approx$ borde + 1,25 MHz), color (video + 3,58 MHz), audio (video + 4,5 MHz).
    \item Diferencia audio/video con \textit{max hold}. Captura.
    \item Si no hay: anotar ``no se dispuso de señal analógica''.
\end{enumerate}
\end{minipage}\hfill
\begin{minipage}[t]{0.49\textwidth}
\textbf{4. TDA}
\begin{enumerate}
    \item Banda UHF completa (fila 6) y captura: los canales digitales son bloques planos de $\approx$ 5,6 MHz.
    \item Anotar cada canal: f central y nivel de meseta.
    \item Elegir 1 o 2 canales (fila 7): ancho entre flancos a $-3$ dB y potencia de canal con 5,57 MHz (si el equipo la tiene; si no, anotar nivel de meseta y RBW). Captura.
    \item Antena al televisor, búsqueda automática, foto de la lista de canales.
    \item Comparar analizador vs. televisor: canales que se ven pero no se decodifican, y al revés.
\end{enumerate}

\vspace{4pt}
\textbf{5. AM con generador (si el docente lo habilita)}
\begin{enumerate}
    \item AM con un tono de 1 kHz e índice cercano a 1.
    \item Medir nivel de portadora y de una banda lateral.
    \item $m = 2 \cdot 10^{\Delta dB/20}$. Con $m = 1$ la banda lateral queda $\approx -6$ dBc.
\end{enumerate}

\vspace{4pt}
\textbf{Recordatorios}
\begin{itemize}
    \item Fotografiar también la pantalla del televisor con los canales.
    \item Anotar cualquier dificultad (falta de señal, interferencias, limitaciones del equipo): va en las conclusiones.
    \item Las capturas van en \archivo{Figuras/mediciones/} con los nombres del cuadro.
\end{itemize}
\end{minipage}

\newpage
\begin{landscape}
{\Large\textbf{Planilla de registro}}\\[6pt]
\renewcommand{\arraystretch}{1.85}
\begin{tabular}{|L{3.0cm}|L{2.4cm}|L{2.0cm}|L{2.2cm}|L{1.6cm}|L{1.5cm}|L{1.5cm}|L{1.6cm}|L{6.8cm}|}
    \hline
    \textbf{Medición} & \textbf{f [kHz / MHz]} & \textbf{Nivel [dBm]} & \textbf{Ancho [kHz / MHz]} & \textbf{Ref. lvl} & \textbf{RBW} & \textbf{VBW} & \textbf{Aten.} & \textbf{Observaciones} \\
    \hline
    AM banda & & & -- & & & & & \\ \hline
    AM emisora & & & & & & & & Emisora: \\ \hline
    FM banda & & & -- & & & & & \\ \hline
    FM emisora & & & & & & & & Emisora / canal: \\ \hline
    FM vecinas & & & & & & & & \\ \hline
    TV analógica & & & & & & & & Canal: \quad A/V [dB]: \\ \hline
    TDA banda & -- & & -- & & & & & Canales vistos: \\ \hline
    TDA canal 1 & & & & & & & & Canal: \quad Pot. canal: \\ \hline
    TDA canal 2 & & & & & & & & Canal: \quad Pot. canal: \\ \hline
    TDA canal 3 & & & & & & & & Canal: \quad Pot. canal: \\ \hline
    AM generador & & & & & & & & Portadora: \quad BL: \quad $m$: \\ \hline
\end{tabular}

\vspace{10pt}
\renewcommand{\arraystretch}{1.5}
\begin{tabular}{|L{4cm}|L{6.5cm}|L{3cm}|L{9.2cm}|}
    \hline
    \multicolumn{4}{|l|}{\textbf{Búsqueda de canales en el televisor}} \\ \hline
    \textbf{Canal físico} & \textbf{Servicios encontrados} & \textbf{¿Se ve bien?} & \textbf{Observaciones} \\ \hline
    & & & \\ \hline
    & & & \\ \hline
    & & & \\ \hline
    & & & \\ \hline
    & & & \\ \hline
\end{tabular}

\vspace{8pt}
\textbf{Equipos:} Analizador \underline{\hspace{4.5cm}} \quad Antena AM \underline{\hspace{3.5cm}} \quad Antena FM/TV \underline{\hspace{3.5cm}} \quad Televisor \underline{\hspace{3.5cm}}
\end{landscape}

\end{document}
